
\chapter[关于 PGFPlots：预备知识]{关于 {\normalfont\PGFPlots{}}：预备知识}

本节包含关于升级、开发团队、安装（如果你需要手动安装的话）以及
故障排查的信息。除了升级说明之外，你可以完全跳过本节。

\PGFPlots{} 完全构建在 \Tikz{}/\PGF{} 之上。了解 \Tikz{} 会简化使用
\PGFPlots{} 的工作，尽管这并非必需。

\section{组成部分}

\PGFPlots{} 由两个组件构成：
%
\begin{enumerate}
    \item 绘图组件（即你当前正在阅读的部分），以及
    \item \PGFPlotstable{} 组件，它简化数字表格的数字格式化
        与后处理。它作为一个独立的
        宏包提供，并拥有自己的手册
        \href{file:pgfplotstable.pdf}{pgfplotstable.pdf}。
\end{enumerate}


\section{升级说明}

本版本提供了大量改进，感兴趣的读者可以在 \texttt{ChangeLog} 中找到
全部细节。不过，以下几项变更值得特别注意。

\PGFPlots{} 中一个共通的设置是键 |compat|：强烈建议始终把它写进你的
|.tex| 文件。虽然这个键给最终用户带来了一些工作，但它也解决了一个常见需求：
它确保你的 |.tex| 文件总是产生相同的输出，即使你安装了新版本的
\PGFPlots{}。另一方面，它让我们作为维护者能够修复软件缺陷并引入行为变更，
前提是这些变更只影响具有合适 |compat| 级别的文档。|compat| 键的确切影响、
可选取值及其含义将在以下各节中说明。


\subsection{新增可选功能}

\PGFPlots{} 在编写时就考虑了向后兼容：旧的 \TeX{}
文件应当无需修改即可编译，外观也不发生变化。不过，新功能偶尔会导致不同的行为。
在这种情况下，\PGFPlots{} 会停用新功能。\footnote{如果出现破坏向后兼容的情况，
我们深表歉意——并请你提交 bug 报告。我们会负责处理。}

任何导致向后兼容问题的新功能或错误修复都需要\emph{手动}且\emph{显式地}
启用。为此，你应当使用
%
\input pgfplots.preliminaries.compatcurrent.tmp

%
\noindent 放在你的导言区中。这将配置兼容层。

你至少应使用 |compat=1.3|。建议的取值会在运行 \TeX{}
之后打印到 |.log| 文件中。

以下是一份变更清单，除非使用合适的
|compat| 级别，否则这些变更处于停用状态：
%
\begin{enumerate}
    \item \PGFPlots{} 1.18 和 1.17 没有额外约束，在
        兼容级别方面与 1.16 相同。
    \item \PGFPlots{} 1.16 没有额外约束，在
        兼容级别方面与 1.15 相同。
    \item \PGFPlots{} 1.15 启用 |3d log sampling|，并修复
        柱状图的 |clip limits| 问题。
    \item \PGFPlots{} 1.14 改变了系统处理非均匀色图的方式，
        并启用高级 |colormap| 操作（参见
        |of colormap|）。
    \item \PGFPlots{} 1.13 修复极坐标轴中的坐标轴标签，并确保
        |shader=flat| 所选颜色与 |z buffer| 和
        |mesh/ordering| 无关。此外，它为所有堆叠柱状图启用
        |stack negative=separate|。保留较旧的兼容级别是为了
        保留最终用户的临时变通方案。
    \item \PGFPlots{} 1.12 启用 |lua backend| 并定义
        |boxplot/estimator=Excel|。
    \item \PGFPlots{} 1.11 改变了 |axis cs|：它现在是默认
        坐标系。如果你写 |\draw (1,2) -- (2,2);|，\PGFPlots{}
        会自动把它视为
        |\draw (axis cs:1,2) -- (axis cs:2,2);|。
    \item \PGFPlots{} 1.10 在兼容性方面与 1.9 没有区别。
    \item \PGFPlots{} 1.9 附带一个预设，用于组合 |ybar stacked| 和
        |nodes near coords|。此外，它会抑制堆叠柱状图中的空增量。
        要启用这个新预设，你必须
        使用 |compat=1.9| 或更高版本。
    \item \PGFPlots{} 1.8 对标签与刻度比例标签的对齐方式进行了
        新的修订。此外，它改进了
        |hide axis| 的包围盒。该修订通过 |compat=1.8| 或
        更高版本启用。

        配置 |compat=1.8| 是\emph{修复}
        三维坐标轴中 |axis lines=center| 的\emph{必要}条件。
    \item \PGFPlots{} 1.7 新增了以坐标轴单位定义柱宽的选项。
        这些通过 |compat=1.7| 或更高版本启用。
    \item \PGFPlots{} 1.6 新增了用于更精确缩放的选项，以及
        针对 |\addplot3 graphics| 的更多缩放选项。这些通过
        |compat=1.6| 或更高版本启用。
    \item \PGFPlots{} 1.5.1 将圆和椭圆的半径解释为
        \PGFPlots{} 坐标（旧版本使用 \pgfname{} 单位向量，
        它们与 \PGFPlots 没有直接关系）。换句话说：从 1.5.1 版开始，
        可以在坐标轴内部编写 |\draw circle[radius=5]|。这需要
        |\pgfplotsset{compat=1.5.1}| 或更高版本。

        若没有此兼容设置，圆和椭圆将像早期版本那样使用
        \pgfname{} 的底层画布单位。
    \item \PGFPlots{} 1.5 默认使用 |log origin=0|（这会影响
        对数柱状图或堆叠对数图）。较旧的版本
        保持 |log origin=infty|。这需要 |\pgfplotsset{compat=1.5}| 或
        更高版本。
    \item \PGFPlots{} 1.4 修复了若干较小的 bug，
        与早期版本相比可能产生约 $1$--$2\text{pt}$ 的差异。
        这需要 |\pgfplotsset{compat=1.4}| 或更高版本。
    \item \PGFPlots{} 1.3 带来了用户界面上的改进。
        旧版本中“行为选项”与“样式
        选项”之间的技术区分不再必要（不过仍
        完全支持）。

        这总是启用的。
    \item \PGFPlots{} 1.3 有一个新功能，允许自动\emph{把坐标轴标签
        紧贴刻度标签移动}。强烈推荐使用。它需要
        |\pgfplotsset{compat=1.3}| 或更高版本。

        由于这会影响间距，默认情况下并不启用。
    \item \PGFPlots{} 1.3 支持反向坐标轴。不再需要
        使用带负单位的变通方案。\pgfkeys{/pdflinks/search key
        prefixes in/.add={/pgfplots/,}{}}

        请查看 |x dir=reverse| 键。

        现有的变通方案仍能正常工作。使用
        |\pgfplotsset{compat=1.3}| 或更高版本，与 |x dir=reverse| 一起
        切换到新版本。
\end{enumerate}


\subsection{可能需要注意的旧功能}

\begin{enumerate}
    \item |scatter/classes| 功能从 1.3 版起会生成恰当的图例。
        这可能会改变已有使用 |scatter/classes| 的图的图例
        外观。
    \item 从 \PGFPlots{} 1.1 开始，\emph{不应再使用} |\tikzstyle|
        来设置 \PGFPlots{} 选项。

        尽管 |\tikzstyle| 仍被某些较旧的 \PGFPlots{}
        选项支持，但你应当把任何出现的 |\tikzstyle| 替换为
        |\pgfplotsset{|\meta{样式名称}|/.style={|\meta{键值列表}|}}|
        或相应的 |/.append style| 变体。更多细节参见
        Section~\ref{sec:styles}。
\end{enumerate}
%
对于这些变更带来的任何不便，我深表歉意。

\begin{pgfplotskey}{%
    compat=\mchoice{
        1.18,
        1.17,
        1.16,
        1.15,
        1.14,
        1.13,
        1.12,
        1.11,
        1.10,
        1.9,
        1.8,
        1.7,
        1.6,
        1.5.1,
        1.5,
        1.4,
        1.3,
        pre 1.3,
        default
    } (initially default)}
    导言区配置
    %
\input pgfplots.preliminaries.compatcurrent.tmp
    %
    允许在向后兼容与最新功能之间进行选择。
    这个键被设计为文档中遇到的第一个 \PGFPlots{} 键，因为它会准备
    全局选项。

    有时，你可能想在同一
    文档中使用不同的版本。这时，提供
    %
\begin{codeexample}[code only]
\begin{figure}
    \pgfplotsset{compat=1.4}
    ...
    \caption{...}
\end{figure}
\end{codeexample}
    %
    \noindent 以便把兼容设置限制在当前
    上下文中（本例中即 |figure| 环境）。

    查看你的 |.log| 文件输出，即可得到 |compat| 的建议取值。

    使用 |\pgfplotsset{compat=default}| 恢复出厂设置。

    尽管通常没有必要，也可以只启用
    选定的变更，而在总体上保持与旧版本的兼容：
    %
    \begin{pgfplotskeylist}{%
        compat/path replacement=\meta{version},
        compat/labels=\meta{version},
        compat/scaling=\meta{version},
        compat/scale mode=\meta{version},
        compat/empty line=\meta{version},
        compat/plot3graphics=\meta{version},
        compat/bar nodes=\meta{version},
        compat/BB=\meta{version},
        compat/bar width by units=\meta{version},
        compat/pgfpoint substitution=\meta{version},
        compat/general=\meta{version}
    }
        假设我们有一个使用 |\pgfplotsset{compat=1.3}|
        的文档，并且你想保持这种方式。

        此外，你发现 1.5.1 版支持圆和
        椭圆。这时，使用
\begin{codeexample}[]
% preamble:
\pgfplotsset{
    compat=1.3,
    compat/path replacement=1.5.1,
}
\begin{tikzpicture}
\begin{axis}[
    extra x ticks={-2,2},
    extra y ticks={-2,2},
    extra tick style={grid=major}]
    \addplot {x};
    \draw (axis cs:0,0) circle[radius=2];
\end{axis}
\end{tikzpicture}
\end{codeexample}

        所有这些键都接受 |compat| 键的可选值。

        |compat/path replacement| 键控制圆和
        椭圆半径的解释方式。

        |compat/labels| 键控制坐标轴标签的对齐方式：要么
        紧邻刻度，要么使用绝对偏移。从 |1.8| 起，它还
        启用了一套全新的坐标轴标签样式修订。在大多数
        情况下，你看不出区别——但它修复了
        三维坐标轴中的 |axis lines=center|。

        |compat/scaling| 键控制 1.4 和 1.6 版中引入的
        一些 bug 修复：为了提高精度，它们可能带来轻微的
        缩放差异。

        |compat/plot3graphics| 控制
        |\addplot3 graphics| 的新功能。

        |compat/scale mode| 允许启用/禁用警告“The
        content of your 3d axis has CHANGED compared to previous versions”，
        因为 |axis equal| 和 |unit vector ratio| 功能在
        1.6 之前的所有版本中都是坏的，并在 1.6 中得到了修复。

        |compat/empty line| 允许在输入文件中写入空行以
        产生一次跳跃。这需要 |compat=1.4| 或更新版本。详见
        |empty line|。

        |compat/BB| 改变包围盒，使其即使在
        |hide axis| 的情况下也紧凑贴合。

        |compat/bar width by units| 允许用坐标轴单位来表示
        |bar width=1|（即以坐标轴单位计）。

        |compat/bar nodes| 启用 |ybar stacked| 和
        |nodes near coords| 的预设。此外，它为
        堆叠柱状图启用 |stacked ignores zero|。

        |compat/general| 键控制 |log origin|、|lua backend|、
        |enable tick line clipping| 和 |boxplot/estimator|。

        |compat/pgfpoint substitution| 键决定 |axis cs| 是否为
        默认坐标系（自 1.11 起）。

        详细效果可参见本节开头。
    \end{pgfplotskeylist}

    取值 \meta{version} 可以是 |default|、一个版本号，或 |newest|。
    取值 |default| 与 |pre 1.3| 相同（除无关紧要的差异外）。
    强烈\emph{不建议}使用 |newest|：它可能根据 \PGFPlots{}
    的当前版本而改变你的文档。请查看你的 |.log| 文件，
    以获得关于最佳版本的建议。
\end{pgfplotskey}


\section{开发团队}

\PGFPlots{} 主要由 Christian Feuersänger 编写，Pascal Wolkotte 和
Nick Papior Andersen 作为业余项目贡献了许多改进。我们希望它有用，
并能提供有价值的图。

如果你有兴趣编写一些东西但不知从何入手，可以参阅
随 \PGFPlots{} 提供的辅助手册
\href{file:TeX-programming-notes.pdf}{TeX-programming-notes.pdf}。它
远非完整，但也许是一个好的
起点（至少能引导你找到更多文献）。


\section{致谢}

感谢上帝赐予我享受编程的无数时光。我感谢 Pascal Wolkotte 和
Nick Papior Andersen 作为开发团队一员所付出的编程努力与贡献。我感谢 Francesco Poli 为 |contour lua| 实现 |lua| 算法。
我感谢 Jürnjakob Dugge 贡献了
|hist/density|、用于 \verbpdfref{\addplot3} |graphics| 的 matlab 脚本、出色的
用户论坛帮助以及有价值的 bug 报告。我感谢 Stefan Tibus，他贡献了
|\addplot shell| 功能。我感谢 Tom Cashman 贡献了
|reverse legend| 功能。特别感谢 Stefan Pinnow 持续
测试 \PGFPlots{}、讨论需求、提出功能请求和
bug 修复（从而带来众多质量改进），以及移植并
集成 |colorbrewer| 库。此外，我感谢 Schweitzer 教授
富有成效的讨论，以及他最初鼓励我开始这样一个
宏包。感谢 Meine 博士的想法和建议。特别感谢
Markus Böhning 校对了 \PGF{}、\PGFPlots{} 和
\PGFPlotstable{} 的所有手册。感谢 Vincent A. Traag 把 |colorbrewer| 配色
带给 \PGFPlots{}。也感谢众多国际贡献者，他们
提出了功能请求、发现了 bug，或仅仅是改进了
手册！

最后但同样重要的是，我感谢 Till Tantau 和 Mark Wibrow 打造了出色的
图形（及更多）宏包 \PGF{} 和 \Tikz{}，它们正是
\PGFPlots{} 的基础。


% main=manual.tex

\section{安装与先决条件}

\subsection{许可}

本程序是自由软件：你可以按照自由软件基金会发布的 GNU 通用公共许可证
（GNU General Public License）第 3 版，或（由你选择）任何更新版本的
条款，重新分发和/或修改它。

本程序的分发希望它会有用，但不提供任何担保；甚至不包含
对适销性或特定用途适用性的默示担保。更多细节请参见 GNU 通用公共许可证。

GNU 通用公共许可证的副本可在宏包文件
%
\begin{verbatim}
doc/latex/pgfplots/gpl-3.0.txt
\end{verbatim}
%
中找到。你也可以访问~\url{http://www.gnu.org/licenses}。


\subsection{先决条件}

\PGFPlots{} 需要 \PGF{}。一般来说，你应使用最新稳定
版本的 \PGF{}。\PGFPlots{} 的使用方式是在
%
\input pgfplots.preliminaries.compatcurrent.tmp
%
\noindent 你的导言区中（关于如何在 Con\TeX{}t 和 plain \TeX{} 中使用，
参见 Section~\ref{sec:tex:dialects}）。

应当添加 |compat=|\meta{你的版本} 这一项以启用新功能，
更多细节参见 |compat| 键的文档。


%\subsection{Installation}

有若干种方法可以告诉 \TeX{} 到哪里寻找这些文件。请选择
最符合你需求的选项。


\subsection{在 Windows 上安装}

Windows 用户通常使用 MiK\TeX{}，它会自动下载最新稳定的宏包
版本。在这里你无需手动安装任何东西。

如果你想安装比 MiK\TeX{} 所附带版本更新的 \PGFPlots{}，
可以按如下步骤进行。MiK\TeX{} 提供了在其自身的 \TeX{} 目录结构
（TDS）中本地安装宏包的功能。基本思路是：把 \PGFPlots{} 解压到你选择的
某个目录，并配置 MiK\TeX{} 宏包管理器，让它以高于默认路径的
优先级使用这个特定目录。如果你想这样做，请通过“开始 $\gg$ 程序 $\gg$ MiK\TeX{} $\gg$
Settings”启动 MiK\TeX{} 设置程序。在那里，使用“Roots”菜单页。
它初始配置中包含 MiK\TeX{} 宏包目录。使用“Add”选择解压后的
\PGFPlots{} 所在目录。然后，使用“Up”按钮把新添加的路径
移到列表顶端。再按“Ok”。对于 MiK\TeX{}
2.8，你可能需要取消勾选“Show MiK\TeX{}-maintained root directories”
按钮，才能看到新安装的路径。

如果所提供的目录不符合 TDS 规范（细节参见
Section~\ref{pgfplots:tds}），MiK\TeX{} 会报错，所以你在这里不能提供错误的
目录。这种方法对其他宏包同样有效，但有些宏包在被 MiK\TeX{}
接受之前可能需要先做一些目录结构上的整理。


\subsection{安装 Linux 宏包}

通常，可以使用 \TeX{} 宏包管理器来安装 \PGFPlots{}。一个
常见的发行版是 \TeX{}Live。在这种情况下，你可以输入
%
\begin{codeexample}[code only]
sudo tlmgr install pgfplots
\end{codeexample}
%
\noindent 来安装 \PGFPlots{}。


\subsection{安装到任意目录——\texttt{TEXINPUTS} 变量}

你可以简单地把 \PGFPlots{} 安装到硬盘上的任意位置，例如
安装到
%
\begin{verbatim}
/foo/bar/pgfplots.
\end{verbatim}
%
然后，把 \texttt{TEXINPUTS} 变量设置为
%
\begin{verbatim}
TEXINPUTS=/foo/bar/pgfplots/tex//:
\end{verbatim}
%
末尾的~`\texttt{:}'告诉 \TeX{} 在 \lstinline!/foo/bar/pgfplots!
之后再检查默认搜索路径。双斜杠~`\texttt{//}'告诉 \TeX{}
搜索所有子目录。

如果 \texttt{TEXINPUTS} 变量中已经有内容，你可以把
上面这一行追加到现有的 \texttt{TEXINPUTS} 内容之后。

此外，你还应当设置 |TEXDOCS|，
%
\begin{verbatim}
TEXDOCS=/foo/bar/pgfplots/doc//:
\end{verbatim}
%
以便 \TeX{} 文档系统能找到 |pgfplots.pdf| 和
|pgfplotstable.pdf| 这两个文件（在某些系统上，之后只需使用
|texdoc pgfplots| 即可）。

从 \PGFPlots{} 1.12 开始，你可能还需要设置 \texttt{LUAINPUTS}：
%
\begin{verbatim}
LUAINPUTS=/foo/bar/pgfplots//:
\end{verbatim}
%
通常即可满足需要。

关于如何设置环境变量，请参阅你操作系统的手册。


\subsection{安装到本地符合 TDS 规范的 \texttt{texmf} 目录}
\label{pgfplots:tds}

\PGFPlots{} 采用符合“\TeX{} 目录结构”（TDS）的
目录结构分发，因此你可以简单地把文件解压到一个会被
\TeX{} 自动搜索的目录中。例如，在 Linux 系统上，这样的目录是
|~/texmf|。

把 \PGFPlots{} 复制到某个本地 \texttt{texmf} 目录，例如 \lstinline!~/texmf!。
你至少需要 \PGFPlots{} 的 |tex/generic/pgfplots| 和
|tex/latex/pgfplots| 目录。然后，运行 \lstinline!texhash!（或某个等效的、
特定于你所用 \TeX{} 发行版的路径更新命令）。

TDS 由若干个子目录构成，根据所请求内容的不同，它们会被
分别搜索：子目录
|doc/latex/|\meta{宏包} 用于（\LaTeX{}）文档，
子目录 |doc/generic/|\meta{宏包} 用于同时适用于
\LaTeX{} 和其他 \TeX{} 方言（例如 plain \TeX{} 和 Con\TeX{}t，它们也有
各自对应的子目录）的文档。

类似地，每当请求 \LaTeX{} 宏包时，会搜索
|tex/latex/|\meta{宏包} 子目录。对于同时适用于 \LaTeX{} \emph{以及}
其他 \TeX{} 方言的宏包，会搜索 |tex/generic/|\meta{宏包}
子目录。

别忘了运行 \lstinline!texhash!。

\subsection{其他方法均失败时的安装\ldots}

如果 \TeX{} 仍然找不到你的文件，你可以把所有 \lstinline!.sty! 和
所有 |.code.tex| 文件（也许还包括所有 |.def| 文件）复制到你当前
项目的工作目录中。实际上，你需要
|tex/latex/pgfplots| 和 |tex/generic/pgfplots| 子目录中的
全部内容。

关于宏包安装的更多信息，请参阅 \url{http://www.ctan.org/installationadvice/}。


\section{故障排查——错误信息}

本节讨论使用 \PGFPlots{} 时可能出现的一些问题。
部分错误信息列于索引中，请看本手册末尾
（“Errors”一节）。


\subsection{可用 Dimen 寄存器不足的问题}

为避免 \PGF{} 和 \PGFPlots{} 所需的大量 \TeX{} 寄存器
引发问题，你可能希望把
%
\begin{verbatim}
\usepackage{etex}
\end{verbatim}
%
作为第一个宏包引入。这在大多数情况下可以避免“no room for a new
dimen”\index{Error Messages!No room for a new dimen} 问题。它应当
适用于任何现代的 \TeX{} 安装（它会激活
$\varepsilon$-\TeX{} 扩展）。


\subsection{Dimension Too Large 错误}

\PGF{} 的核心数学引擎依赖 \TeX{} 寄存器来执行
快速算术运算。为计算 $50+299$，它实际上计算的是 |50pt+299pt|，然后
去掉结果的 |pt| 后缀。由于 \TeX{} 寄存器只能容纳
不超过 $\pm 16384$ 的数，一旦结果超出允许范围，就会出现像“Dimension too large”这样的
溢出错误信息。正常情况下，这绝不应该
发生——\PGFPlots{} 使用数据范围为 $\pm 10^{324}$ 的浮点单元，
并自动完成所有映射。不过，在某些情况下这仍会
失败。其中一些情况如下：
%
\begin{enumerate}
    \item 坐标轴范围（例如 $x$ 的范围）变得\emph{相对}
        很小。像 $[10^{-17},10^{-16}]$ 这样绝对很小的范围
        没有问题。但如果你的坐标轴范围是
        $[1.99999999,2]$，即需要大量有效数字，
        这就可能出问题。

        恐怕我在这里帮不上忙：你可能需要在
        \PGFPlots{} 处理数据之前，先以某种方式整理好数据。
    \item 如果你只查看数据范围中非常小的一部分，
        也可能发生这种情况。

        例如，当你的输入数据范围为 $x\in
        [0,10^6]$，而你指定 |xmax=10| 时，就会发生这种情况。

        在这种情况下，请考虑使用
        |restrict x to domain*=|\meta{min}|:|\meta{max}
        键，其中 \meta{min} 和 \meta{max} 应当
        （比方说）为坐标轴上下限的四倍（细节参见
        page~\pageref{key:restrict:x:to:domain}）。
    \item 如果 $x$ 和 $y$ 的尺度相差非常悬殊，
        |axis equal| 键会被弄糊涂。
    \item 你可能发现了一个 bug——请联系开发者。
\end{enumerate}


\subsection{DVI 查看器与 \texttt{dvipdfm} 的限制}
\label{sec:drivers}

\PGF{} 兼容
%
\begin{itemize}
    \item \lstinline!latex!/\lstinline!dvips!,
    \item \lstinline!latex!/\lstinline!dvipdfm!,
    \item \lstinline!pdflatex!,
    \item $\vdots$
\end{itemize}
%
不过，存在一些限制：我所知道的 DVI 查看器都无法
查看 \PGF{} 的输出（因而也无法查看 \PGFPlots{} 的输出）。
毕竟，DVI 从来就不是为绘制文字和水平/垂直直线
以外的东西而设计的。你需要查看 postscript 文件或
PDF 文件。

其次，DVI/PDF 组合并不支持所有类型的着色
（例如，|shader=interp| 仅适用于 |dvips|、|pdftex|、
|dvipdfmx| 和 |xetex| 驱动）。

此外，\PGF{} 需要知道一个\emph{驱动}，以便把 DVI 文件
转换为期望的输出。根据你的系统，你需要以下
选项：
%
\begin{itemize}
    \item \lstinline!latex!/\lstinline!dvips! 不需要任何特殊设置，
        因为 \lstinline!dvips! 是你调用
        \lstinline!latex! 时的默认驱动。
    \item \lstinline!pdflatex! 也可以直接使用，因为
        \lstinline!pdflatex! 会被自动检测到。
    \item \lstinline!lualatex! 同样会被自动检测到。
    \item \lstinline!latex!/\lstinline!dvipdfm! 需要使用
        %
\begin{verbatim}
\def\pgfsysdriver{pgfsys-dvipdfm.def}
%\def\pgfsysdriver{pgfsys-pdftex.def}
%\def\pgfsysdriver{pgfsys-dvips.def}
%\def\pgfsysdriver{pgfsys-dvipdfmx.def}
%\def\pgfsysdriver{pgfsys-xetex.def}
%\def\pgfsysdriver{pgfsys-luatex.def}
\usepackage{pgfplots}.
\end{verbatim}
        %
        被注释掉的命令可用于显式设置其他驱动。
\end{itemize}
%
如果你还有进一步的问题，请阅读~\cite[Sections~7.2.1 and 7.2.2]{tikz}
中相应的章节。这些章节也包含特定驱动的
局限性。

完全不会产生任何问题的选择是 |pdflatex|。


\subsection{\TeX{} 内存容量不足的问题}

\PGFPlots{} 可以处理从小到中等规模的图。不过，\TeX{}
从来就不是为数据图而设计的——你最终会遇到
内存容量不足的问题。如何扩大内存，参见
Section~\ref{sec:pgfplots:optimization}。


\subsection{语言设置与活动字符的问题}

\PGF{} 和 \PGFPlots{} 都使用大量活动字符——这可能会
与其他定义活动字符的宏包产生
不兼容。兼容性比早期版本有所改善，但仍可能
是个问题。本手册分别使用英文和法文的 |babel| 宏包编译，
|german| 宏包同样可行。如果你遇到任何问题，请告诉我。
有时，临时禁用活动字符可能会奏效（|babel|
提供了这样的命令）。


\subsection{其他问题}

请查阅 \url{https://github.com/pgf-tikz/pgfplots/issues} 上的 issue 列表
或邮件列表 \url{https://tug.org/mailman/listinfo/pgf-tikz}。也许
之前也有人遇到过你的问题，并且他可能已经想出了
一个解决办法。

如果你遇到的是一个不同的问题，请新开一个 issue。Matrix 网络上还有一个聊天室，
地址是 \href{https://matrix.to/#/#pgf-tikz:matrix.org}{\texttt{\#pgf-tikz:matrix.org}}。

